\chapter{Anemia}
\index{Anemia}

\section*{Definition and Overview}
Anemia is a condition in which the body does not have enough healthy red blood cells, or enough of the iron-containing protein hemoglobin, to carry oxygen around the body. Because oxygen is needed to produce energy, people with anemia often feel tired and weak.

Anemia is not a single disease but a sign of an underlying problem, and there are many types, including iron-deficiency anemia, vitamin B12 and folate deficiency anemia, and anemia caused by chronic disease or blood loss. The cause determines how it is treated, so identifying the type is the first step in care.

\section*{Epidemiology}
Anemia is the most common blood disorder in the world, affecting roughly one in four people globally. It is especially common in women of childbearing age, mainly because of menstrual blood loss, and in children and adolescents whose growing bodies need extra iron.

Iron deficiency is the leading cause worldwide, particularly in low-income countries, while in developed countries it is often linked to heavy periods, pregnancy, poor diet, or hidden bleeding from the stomach or bowel. Anemia of chronic disease is common in older people with conditions such as kidney disease, arthritis, or cancer.

\section*{Aetiology and Pathophysiology}
\begin{itemize}
    \item \textbf{Blood loss:} heavy menstrual periods, bleeding from the stomach or bowel, such as from ulcers, polyps, or bowel cancer, or injury deplete the body's iron stores
    \item \textbf{Reduced production:} the bone marrow needs iron, vitamin B12, folate, and other nutrients to make red blood cells, and deficiency of any of these slows production
    \item \textbf{Chronic disease:} long-term illness, inflammation, kidney disease, or cancer can reduce the marrow's ability to use iron, even when iron levels are normal
    \item \textbf{Increased destruction:} in hemolytic anemias, red blood cells are destroyed faster than they can be replaced
    \item \textbf{Genetic conditions:} sickle cell disease, thalassemia, and other inherited anemias affect red cell shape or production
    \item \textbf{Mechanism:} fewer or smaller red blood cells, or less hemoglobin within them, means less oxygen is delivered to the tissues, producing symptoms
\end{itemize}

\section*{Clinical Features}
\begin{itemize}
    \item Fatigue, weakness, and a general lack of energy
    \item Pallor, with paler skin than usual, and paleness of the inside of the lower eyelids and nail beds
    \item Shortness of breath on exertion
    \item Dizziness, light-headedness, or headaches
    \item Palpitations, or a fast or pounding heartbeat
    \item Cold hands and feet
    \item Less common features: brittle or spoon-shaped nails, a sore tongue, unusual cravings for non-food items (pica) such as ice or dirt, and, in severe cases, chest pain or fainting
\end{itemize}

\section*{Differential Diagnosis}
\begin{itemize}
    \item Iron-deficiency anemia, the most common type, from blood loss or poor iron intake
    \item Vitamin B12 or folate deficiency anemia, which can also cause tingling in the hands and feet
    \item Anemia of chronic disease, from ongoing inflammation or organ failure
    \item Hemolytic anemia, from rapid destruction of red blood cells
    \item Bone marrow disorders, such as aplastic anemia or leukemia
    \item Genetic disorders such as the thalassemia trait, which can resemble iron deficiency
    \item Conditions that mimic the symptoms of anemia, such as an underactive thyroid, depression, or chronic fatigue
\end{itemize}

\section*{When to Seek Medical Care}
See a doctor if you have persistent tiredness, paleness, shortness of breath, or palpitations, or if you have been losing more blood than usual. Anemia can have serious underlying causes, so it should always be investigated.

\textbf{Contact your local emergency services for:}
\begin{itemize}
    \item Chest pain, severe breathlessness, or fainting
    \item Very heavy bleeding that you cannot control
    \item Sudden severe weakness or confusion
    \item Coughing up blood, or passing large amounts of blood in the stools or vomit
\end{itemize}

\section*{Diagnosis and Investigations}
\begin{itemize}
    \item \textbf{History:} diet, menstrual and pregnancy history, medications, family history, and any symptoms of bleeding
    \item \textbf{Full blood count:} measures red blood cells, hemoglobin, and related indices, and confirms whether anemia is present and how severe it is
    \item \textbf{Iron studies:} blood tests for iron, ferritin, and related markers to diagnose iron deficiency
    \item \textbf{Vitamin B12 and folate levels:} to detect deficiency
    \item \textbf{Blood film:} examination of the cells under a microscope can point to the type of anemia
    \item \textbf{Hemolysis tests:} for anemia caused by red cell destruction
    \item \textbf{Further tests depending on the cause:} endoscopy or colonoscopy for suspected bowel bleeding, or assessment of the kidneys and other organs
    \item \textbf{Referral to a hematologist:} when the cause is unclear or the anemia is severe
\end{itemize}

\section*{Management}
\begin{itemize}
    \item \textbf{Treat the cause:} stopping the source of bleeding, treating an underlying disease, or correcting an autoimmune process
    \item \textbf{Iron supplementation:} oral iron tablets are the main treatment for iron-deficiency anemia, sometimes given intravenously if tablets are not tolerated or absorption is poor
    \item \textbf{Vitamin supplements:} vitamin B12, sometimes by injection, and folate for deficiency anemias
    \item \textbf{Diet:} iron-rich foods such as lean red meat, legumes, and leafy greens, together with vitamin C, which helps iron absorption
    \item \textbf{Medicines:} treatments for the underlying condition, such as erythropoiesis-stimulating agents in kidney disease
    \item \textbf{Blood transfusion:} for severe anemia or when rapid correction is needed
    \item \textbf{Follow-up:} repeat blood tests to confirm that blood counts return to normal and that the cause is controlled
\end{itemize}

\section*{Complications}
\begin{itemize}
    \item Severe fatigue that interferes with work, study, and daily life
    \item Reduced exercise tolerance and physical capacity
    \item Heart problems, including palpitations, an enlarged heart, or heart failure, in prolonged severe anemia
    \item Complications in pregnancy, including low birth weight and premature birth
    \item Delayed growth and development in children
    \item A missed or delayed diagnosis of the underlying cause, such as bowel cancer
\end{itemize}

\section*{Prognosis and Outcomes}
The outlook for anemia is generally very good once the cause is identified and corrected. Iron-deficiency anemia usually improves within several weeks of starting treatment, although it may take months to rebuild iron stores. Anemia of chronic disease tends to persist while the underlying illness persists and is managed rather than cured. Some inherited anemias require lifelong care. Most people feel significantly better, with more energy and less breathlessness, once their blood count returns toward normal.

\section*{Prevention and Self-Care}
\begin{itemize}
    \item Eat a diet with enough iron, vitamin B12, and folate, including lean meat, fish, eggs, legumes, and green leafy vegetables
    \item Combine iron-rich foods with vitamin C, such as fruit or juice, and avoid tea or coffee with meals, which reduce iron absorption
    \item Take iron supplements during pregnancy as advised by your doctor or midwife
    \item Manage heavy periods with medical advice rather than ignoring the blood loss
    \item Report any bleeding from the stomach or bowel, such as black stools, to a doctor promptly
    \item Attend regular check-ups and repeat blood tests when advised to ensure anemia does not return
\end{itemize}

\section*{Sources and Further Reading}
The following resources provide reliable, peer-reviewed and patient-orientated information on anemia.

\begin{itemize}
    \item NHS. \textit{Overview of anemia, its types, and treatment.} https://www.nhs.uk/conditions/
    \item Mayo Clinic. \textit{Guide to the symptoms, causes, and diagnosis of anemia.} https://www.mayoclinic.org/
    \item Centers for Disease Control and Prevention. \textit{Information on blood disorders and iron deficiency.} https://www.cdc.gov/
    \item MSD Manual. \textit{Professional and patient information on the classification and causes of anemia.} https://www.msdmanuals.com/
    \item MedlinePlus. \textit{Health information on anemia.} https://medlineplus.gov/
    \item NICE. \textit{Clinical guidance on the assessment and management of anemia.} https://www.nice.org.uk/
\end{itemize}
